\documentclass[11pt]{article}
\usepackage[utf8]{inputenc}
\usepackage{amsmath,amssymb}
\usepackage{authblk}
\usepackage[margin=1in]{geometry}
\usepackage[hidelinks]{hyperref}
\usepackage{xurl}
\usepackage{setspace}

\title{Real Physics without Philosophy\\
\large Ontology, Prediction, and the End of Interpretation}
\author[1]{Claes Johnson}
\affil[1]{KTH Royal Institute of Technology, SE-100 44 Stockholm, Sweden \\
\texttt{claesjohnson@gmail.com}}
\date{\today \\[0.5em]
\small\textit{Written in respectful cooperation between Claes Johnson and Claude (Anthropic), with Claude
contributing a balanced, neutral view.}}

\begin{document}
\maketitle

\begin{abstract}
A physical theory is an account of what exists and how it behaves; prediction is derivative --- the behaviour
of what exists, computed forward. By this measure standard quantum mechanics is not, or not yet, real physics:
it is an \emph{epistemology}, a calculus for the outcomes of specific experiments, that declines to say what
is physically present between them. This is not a stance one may take without cost. Prediction
\emph{presupposes} ontology --- a prediction is a statement about something real --- and standard QM meets the
demand only by smuggling in a classical measuring apparatus with an ontology of its own while denying ontology
to the microscopic: the Heisenberg cut, which the theory cannot locate and which renders it parasitic on a
classical world it cannot itself ground. The measurement problem, collapse, and the nonlocality puzzles are
then not deep discoveries but symptoms of a single absence, and they trace to a single object: the wave
function over $3N$-dimensional configuration space, which is not a thing in space and so has no ontology. The
century of ``interpretations'' of quantum mechanics is philosophy hired to supply the ontology the formalism
lacks. RealQM removes the need: it \emph{is} an ontology --- $N$ charge densities in ordinary
three-dimensional space --- from which real physics is computed (atoms, molecules, nuclei), not the outcomes
of measurements but the things themselves. A theory with a clear ontology is simply physics, and needs no
interpretation. Philosophy of physics, on this reading, is the discipline of coping with theories that lack an
ontology; give physics back its ontology, and the discipline has nothing left to do.
\end{abstract}

\setstretch{1.06}

\section{Physics is ontology; prediction is derivative}

A physical theory tells us what the world is made of and how it behaves; from that, what we will see follows.
Newtonian mechanics posits masses at positions with momenta, evolving under forces, and its predictions are
the computed trajectories of those real things. Maxwell's theory posits fields in space, and its predictions
are what those fields do. Continuum mechanics posits densities, stresses and velocities in space, and predicts
their evolution. In every case the prediction is \emph{derivative}: it is the forward computation of the
behaviour of something the theory holds to exist. Remove the ontology --- the ``what exists'' --- and there is
nothing whose behaviour is being computed. Physics without ontology is not a leaner physics; it is not
physics. It is, at most, a rule for anticipating readings, which is a different and lesser thing, and which ---
as we now argue --- cannot even stand on its own.

\section{Prediction presupposes ontology: the Heisenberg cut}

It is often held that quantum mechanics is the exception --- that it predicts the outcomes of measurements and
need commit to nothing further. But a prediction is a statement \emph{about something}: to predict an outcome
is to predict that some real thing will be found in some real state. Standard QM's ``outcomes'' are readings
of an apparatus --- a pointer, a click, a track in an emulsion --- which are macroscopic, classical, and
taken as real. So the theory does have an ontology after all: the classical apparatus. What it refuses is an
ontology for the microscopic system, which it represents only by an amplitude.

The result is the \emph{Heisenberg cut}: a line, nowhere fixed by the theory, above which things are real
(apparatus, outcomes) and below which they are only amplitudes (the system). Two failures follow. First, the
theory cannot say where the cut falls, because nothing in the formalism marks it --- the same electron is a
``system'' in one description and part of an ``apparatus'' in another. Second, and worse, the theory grants to
the macroscopic the very reality it denies to the microscopic, although the macroscopic is \emph{made of} the
microscopic: the pointer is a swarm of the same electrons and nuclei the theory says have no definite state.
This is not a subtlety to be managed but an incoherence. Standard QM's celebrated predictive record is a
record of predicting apparatus readings, and it borrows the reality of those readings from a classical world
it cannot itself provide and cannot consistently withhold from the parts. Prediction without ontology is thus
not economical; it is contradictory --- it must help itself, at the macroscopic end, to exactly the ontology
it refuses at the microscopic.

\section{The measurement problem is the missing ontology}

Seen this way, the measurement problem is not a deep feature awaiting a clever resolution; it is the direct
expression of the missing ontology. Ask \emph{what is measured}, and the theory has no answer, because the
object it calls fundamental --- the amplitude --- is not a thing that could be measured. ``Collapse'' is the
name for the step in which the theory must, without warrant from its own dynamics, convert the amplitude into a
definite fact at the moment a measurement is declared. Schr\"odinger's cat is the reductio of pretending the
amplitude was the reality all along. Each is a symptom of one absence: a theory that said what exists would
say what is measured, and would not need to collapse anything, because the definite fact would have been there
before the measurement, as its object.

\section{One root: the wave function over configuration space}

The absence has a single source. The fundamental object of standard QM, $\Psi(x_1,\dots,x_N)$, lives on a
$3N$-dimensional \emph{configuration space}, not in the three-dimensional space that we, our laboratories, and
our instruments occupy. Kohn stated the working consequence in accepting the Nobel Prize: the many-electron
wave function is not a legitimate scientific concept beyond a handful of electrons, being unrecordable even in
principle. A function on a space that is not physical space is not a thing in the world; it can carry a
probability distribution, but it cannot be an ontology. From this one fact the familiar troubles descend --- the
\emph{arena} problem (in what space does the state exist?), the \emph{individuality} problem (identical
particles with no separate existence), and the nonlocality of a correlated whole that has no parts in space to
be local to. The puzzles of quantum foundations are not many independent mysteries; they are one, the
absent ontology of an object in the wrong space, wearing several faces.

\section{Interpretation as symptom}

The century of ``interpretations'' of quantum mechanics is the diagnostic. Copenhagen declares that the
microscopic has no properties until measured --- it accepts the absence and forbids the question. Many-worlds
keeps the amplitude as real and multiplies the world to fit it. De Broglie--Bohm restores particles with
definite positions beneath the amplitude --- it \emph{adds} an ontology. Objective-collapse models change the
dynamics so the amplitude localises into one. However they differ, each is an attempt either to supply the
ontology the formalism lacks or to legislate that none is needed. That so many mutually incompatible
interpretations coexist, none forced by the equations, is not the mark of a rich theory but of an
\emph{underdetermined} one: the formalism does not say what exists, so the saying is handed to philosophy. A
theory that said what exists would admit no interpretations, for there would be nothing left to interpret.
Philosophy of physics, in its dominant modern form, simply \emph{is} the labour of coping with a physics that
withheld its own ontology --- a discipline whose necessity is a confession.

\section{RealQM: an ontology that predicts real physics}

RealQM removes the need at the root. Its fundamental object is not an amplitude on configuration space but $N$
charge densities $\psi_i^2$, each carried on its own region $\Omega_i$ of ordinary three-dimensional space,
the regions non-overlapping~\cite{Ontology}. This is an ontology in the plain sense: it says what exists ---
extended charge in space --- and where. It is the same \emph{kind} of object that macroscopic continuum
mechanics works with; an atom is a continuum-mechanics problem an \aa ngstr\"om across, not a categorically
different thing. From it real physics is computed --- the total energies of atoms across the periodic table,
covalent and hydrogen bonds, and, reading the nucleus as protons and electrons bound by Coulomb alone, nuclear
binding energies~\cite{RealQM} --- and what is computed is not the reading of an instrument but the thing
itself: the charge distribution, its energy, its geometry.

The foundational puzzles do not arise, because their common source is gone. There is no Heisenberg cut,
because everything is real at every scale --- the apparatus and the system are made of the same charge in the
same space. There is no collapse, because the evolution is deterministic and the density is definite at all
times; ``measurement'' is the ordinary interaction of one real charge distribution with another. Schr\"odinger's
cat is a definite configuration throughout. The apparent randomness of quantum experiments is the epistemic
randomness of initial configurations we neither see nor control, as in any deterministic account of a
complicated system --- ours, not the world's.

We do not claim RealQM is complete. It does not yet deliver the spin-exchange phenomena (the singlet--triplet
splitting, collective magnetism), and the boundary of its reach is still being mapped. But its incompleteness
is the \emph{ordinary} incompleteness of a physical theory --- there is more to compute, and the claims can be
checked and found wanting --- not the structural incoherence of a theory that cannot say what it is about. A
theory that can be \emph{wrong} is doing physics; a theory that cannot say what it is about is doing something
else. RealQM can be wrong, and where it is right it is right about \emph{things}, not about readings. That is
the difference in kind.

\section{The wager, and its reach}

It is worth setting down, plainly and as a conditional, exactly what RealQM claims and how far the claim
extends. The claim is a wager about ontology: that the physics of matter is, at bottom, \emph{real charge in
ordinary three-dimensional space}. If that is so, then RealQM --- the dynamics of exactly that --- can capture
it. The syllogism is valid in form; its force rests on two loads, and honesty lives in both.

\emph{First: how much of physics is real charge?} Where the answer is ``essentially all of it'' --- the
Coulomb-structured world of atoms, molecules, bonds, the periodic table, and a growing list of magnetic and
spectroscopic observables --- RealQM captures the physics, and often more cleanly than the configuration-space
formalism, because that domain simply \emph{is} the electrostatics of real charge. But some of physics is not
manifestly charge alone: light is the electromagnetic \emph{field}, a second real entity, so radiation lives in
a \emph{charge-plus-field} coupling; the statistics of measurement across incompatible settings (the Bell
correlations) are informational as much as material; gravity and the precision relativistic sector stand outside
the present theory. Where the premise is incomplete, so is the reach.

\emph{Second: even granting the ontology, is the dynamics right?} Possessing the right stuff --- charge in real
space --- does not by itself fix the right law. RealQM's particular functional (Coulomb, no self-interaction,
non-overlapping domains, kinetic confinement) is itself a hypothesis, confirmed by its hits (total energies,
covalent and hydrogen bonds, helium's diamagnetic susceptibility) and corrected by its misses (many-electron
valence sizes, exchange pairing). A true ontology still owes a true equation.

So the honest thesis is neither ``RealQM explains everything'' nor ``some residue defeats it,'' but a conditional
under adjudication: \textbf{where physics is real charge in three-dimensional space, RealQM captures it; the open
frontier is to measure how far that reaches, phenomenon by phenomenon} --- extending the picture where charge
alone is not enough (adding the field), and marking honestly where it meets a genuine limit. That discipline ---
capture, extend, or concede, case by case --- is what makes it physics rather than doctrine.

\section{Real physics without philosophy}

The conclusion is in the title. Physics with a clear ontology is simply physics: it says what exists, computes
what it does, and is tested against the world. It needs no interpretation, because it has already said what it
is about --- and one does not interpret a bridge calculation or a weather model, one checks it. The elaborate
philosophy that surrounds quantum mechanics --- the interpretations, the measurement debates, the metaphysics
of the wave function --- is not intrinsic to the physics of the small. It is the price of a formalism that
left three-dimensional space and, with it, its ontology. Return the theory to real space, let the charge
densities be the things they plainly appear to be, and the philosophy is not refuted so much as
\emph{retired}: there is nothing left for it to do. The mess was never in nature; it was in the arena. Change
the arena back, and physics is once again about what there is.

\begin{thebibliography}{9}
\bibitem{RealQM} C.~Johnson, \emph{Quantum Mechanics as Multiphase 3D Continuum Mechanics}, manuscript,
2026; \url{https://claes542.github.io/RealMolecule/gallery.html}.
\bibitem{Ontology} C.~Johnson, \emph{Charge Densities in Real Space: RealQM and the Realist Quest Carried by
Schr\"odinger}, manuscript, 2026.
\bibitem{Born26} M.~Born, \emph{Zur Quantenmechanik der Sto\ss vorg\"ange}, Z.\ Phys.\ \textbf{37}, 863
(1926).
\bibitem{Bell87} J.~S.~Bell, \emph{Speakable and Unspeakable in Quantum Mechanics}, Cambridge University
Press, 1987.
\end{thebibliography}

\end{document}
